\section{安全与观测深讲：没有证据就没有边界}

安全和可观测性必须一起讲。权限检查失败时，要能证明没有改变目标对象；进程被 seccomp 拒绝时，要能看到拒绝的系统调用；cgroup 限制触发时，要能解释是内存、CPU 还是进程数限制；panic 发生时，要保留足够证据。没有观测能力，安全边界只能靠想象。

\subsection{威胁模型写法}

安全分析应先写威胁模型，而不是先列工具名。

\begin{lstlisting}
Threat model:
  attacker controls:
    HTTP request body and headers
    uploaded file content
  attacker does not control:
    service binary
    kernel code
    host root credentials
  assets:
    user database
    service token
    host filesystem
  required boundaries:
    process cannot read host secrets
    process cannot mount or ptrace
    process cannot consume unlimited memory
\end{lstlisting}

这样才能判断 namespace、cgroup、capability、seccomp 和文件权限各自解决哪一部分问题。否则“用了容器”不等于安全，“用了 seccomp”也不等于安全。

\subsection{最小权限的阶段化}

服务启动阶段和运行阶段需要的权限不同。启动可能需要读取配置、绑定端口、打开日志；运行时只需要处理已有 fd 和少量系统调用。权限应随阶段下降。

\begin{lstlisting}
startup:
  read config
  bind port 80
  open log file

drop:
  chroot or mount namespace
  drop capabilities
  set no_new_privs
  install seccomp allowlist
  enter cgroup limits

serve:
  accept, read, write, epoll, futex, clock
\end{lstlisting}

权限下降要可测试。测试不应只证明正常请求成功，还要证明 \code{mount}、\code{ptrace}、打开宿主敏感路径、超额 fork 和超额内存分配都会失败。

\subsection{错误路径的证据}

错误路径也要写日志和计数器，但不能泄露敏感信息。比如权限拒绝应记录主体、对象类型、操作和拒绝原因，但不应把密钥内容写入日志。panic 记录寄存器、栈、锁和任务状态，但不能依赖复杂文件系统路径。

\begin{longtable}{p{0.22\linewidth}p{0.34\linewidth}p{0.32\linewidth}}
\toprule
事件 & 必要证据 & 避免 \\
\midrule
权限拒绝 & uid、capability、对象类型、操作 & 打印敏感文件内容 \\
OOM & cgroup、进程、内存统计、victim & 只写 killed without reason \\
seccomp 拒绝 & syscall number、action、pid & 静默杀进程无诊断 \\
panic & 原因、CPU、寄存器、栈、锁 & panic 路径做复杂阻塞 I/O \\
\bottomrule
\end{longtable}

\subsection{从现象到证据链}

一个好的 OS 调试结论要形成证据链。

\begin{lstlisting}
Symptom:
  p99 latency increased from 50ms to 2s

Evidence chain:
  application trace: requests stuck in write phase
  syscall trace: fsync takes 1.8s
  kernel metrics: dirty pages high
  block trace: queue depth high, flush waits long
  conclusion: tail latency dominated by storage flush
\end{lstlisting}

这种报告比“怀疑磁盘慢”有用，因为它说明了每一层的证据和排除方向。若没有 syscall trace，也可能是应用锁；若没有 block trace，也可能是文件系统 journal；证据链让下一步实验明确。

\subsection{安全与观测练习}

\begin{enumerate}
\tightlist
\item 写一个网络服务的威胁模型，列出攻击者能控制和不能控制的对象。
\item 设计 seccomp allowlist，并说明正常路径需要哪些系统调用。
\item 构造权限拒绝测试，证明失败时目标文件没有被创建或截断。
\item 写一份 OOM 诊断报告，包含 cgroup、进程和内存统计。
\item 给一次 panic 设计最小 crash dump 字段。
\end{enumerate}

\subsection{安全边界的测试矩阵}

安全机制必须能被测试。下面是一组最小矩阵：

\begin{longtable}{p{0.24\linewidth}p{0.34\linewidth}p{0.3\linewidth}}
\toprule
边界 & 测试 & 期望 \\
\midrule
页表隔离 & 用户态读内核地址 & page fault 或进程终止 \\
文件权限 & 无权限 open/write & 返回拒绝且目标不变 \\
capability & drop 后执行特权操作 & 返回 \code{EPERM} \\
namespace & 访问不可见路径或进程 & 不可见或返回不存在 \\
cgroup & 超过内存或进程限制 & OOM/失败限制在 cgroup 内 \\
seccomp & 调用禁止 syscall & 按策略返回 errno 或杀进程 \\
观测 & 触发拒绝和 OOM & 日志/计数器能解释原因 \\
\bottomrule
\end{longtable}

测试安全时要检查“没有发生什么”。权限拒绝后文件不能被截断；seccomp 拒绝后进程不能完成危险操作；cgroup OOM 不能扩大成宿主全局 OOM。负面保证是安全测试的核心。

\subsection{可观测性常见误区}

\begin{warningbox}
\begin{itemize}
\tightlist
\item CPU 低不代表系统没压力，可能都在 off-CPU 等 I/O 或锁。
\item RSS 增长不一定是泄漏，可能是 page cache、arena 或 mmap。
\item p99 变差不一定是平均路径慢，可能是少数 fsync、reclaim 或锁等待。
\item 日志多不等于可观测，结构化字段和请求身份更重要。
\item 没有反证的结论只是猜测。
\end{itemize}
\end{warningbox}

\subsection{本章小结}

\begin{itemize}
\tightlist
\item 安全需要威胁模型，不能只列机制名。
\item 最小权限要随阶段下降，并用测试证明危险操作失败。
\item 错误路径必须留下足够证据，但不能泄露敏感数据。
\item 调试报告要形成现象、证据、反证、结论和边界的链条。
\item 安全和可观测性互相支撑：看不见的边界很难被信任。
\end{itemize}

\subsection{访问控制：主体、客体、操作和上下文}

权限判断可以抽象成四元组：主体是谁，客体是什么，操作是什么，上下文是什么。主体可能是 uid、gid、capability、进程安全标签或容器身份；客体可能是 inode、socket、进程、设备或命名空间对象；操作是 read、write、execute、mount、ptrace 等；上下文包括路径、挂载选项、是否 setuid、是否处于用户命名空间等。

\begin{lstlisting}
may_access(subject, object, operation, context):
  if object.type == inode:
    if dac_denies(subject.uid, object.mode, operation):
      if !has_capability(subject, override_cap(operation)):
        return DENY
  if lsm_policy_denies(subject.label, object.label, operation):
    return DENY
  if mount_flags_deny(context.mount, operation):
    return DENY
  return ALLOW
\end{lstlisting}

真实系统通常叠加多层策略：传统 UNIX DAC、capability、ACL、LSM、namespace 边界和挂载选项。安全 bug 常来自只检查其中一层，或在路径解析和最终对象使用之间发生 TOCTOU。正确做法是尽量在持有对象引用后检查权限，并让检查结果绑定到实际对象而不是字符串路径。

\subsection{capability：把 root 权限拆小}

传统 root 权限过大。capability 把特权拆成多个位，例如绑定低端口、修改网络配置、加载模块、改变文件所有者、绕过 DAC 等。服务启动时可能短暂需要某些 capability，运行时应丢弃。

\begin{longtable}{p{0.28\linewidth}p{0.34\linewidth}p{0.26\linewidth}}
\toprule
能力类型 & 能做什么 & 风险 \\
\midrule
\code{CAP\_NET\_BIND\_SERVICE} & 绑定低端口 & 可在启动后丢弃 \\
\code{CAP\_SYS\_ADMIN} & 大量系统管理操作 & 过大，接近半个 root \\
\code{CAP\_SYS\_PTRACE} & 调试或观察其他进程 & 可读敏感内存 \\
\code{CAP\_DAC\_OVERRIDE} & 绕过文件权限检查 & 可读写本不该访问的文件 \\
\code{CAP\_NET\_ADMIN} & 改网络接口和路由 & 可影响宿主网络 \\
\bottomrule
\end{longtable}

最小权限不是“非 root 运行”一句话。一个非 root 进程若仍有危险 capability，仍可能扩大影响面。反过来，一个 root 启动的进程若完成初始化后清空 capability、设置 \code{no\_new\_privs} 并进入受限 namespace，风险会显著降低。

\subsection{namespace：隔离名字，不隔离内核}

namespace 改变进程看到的名字空间：mount namespace 改变路径树，pid namespace 改变进程编号视图，net namespace 改变网络设备和路由视图，uts namespace 改变 hostname，ipc namespace 改变 System V/POSIX IPC 视图，user namespace 改变 uid/gid 映射。

\begin{longtable}{p{0.22\linewidth}p{0.36\linewidth}p{0.3\linewidth}}
\toprule
namespace & 隔离内容 & 仍然共享 \\
\midrule
mount & 文件系统挂载视图 & 同一内核 VFS 和底层设备风险 \\
pid & 进程编号和父子视图 & 调度器、内核内存、系统调用实现 \\
net & 网卡、路由、端口空间 & 协议栈代码和内核漏洞面 \\
user & uid/gid 映射和部分能力 & 内核仍必须正确限制能力转换 \\
ipc & IPC 对象名字空间 & 内核 IPC 实现 \\
\bottomrule
\end{longtable}

容器不是虚拟机，因为容器共享宿主内核。namespace 隔离的是名字和视图，不是内核代码本身。若内核某系统调用存在提权漏洞，容器内进程可能利用同一个内核漏洞逃逸。因此容器安全通常要叠加 seccomp、capability、只读挂载、cgroup 和及时更新内核。

\subsection{cgroup：资源边界也是安全边界}

没有资源限制，攻击者可以通过 fork bomb、内存耗尽、磁盘写满或 CPU 忙等影响整机。cgroup 给进程集合设置资源边界，例如 CPU 权重/配额、内存上限、进程数、I/O 权重和设备访问。

\begin{longtable}{p{0.22\linewidth}p{0.34\linewidth}p{0.32\linewidth}}
\toprule
资源 & 没有限制的风险 & 期望证据 \\
\midrule
CPU & 忙等拖慢其他服务 & throttled time、quota hit 计数 \\
memory & OOM 扩散到宿主 & cgroup OOM 事件、victim、统计 \\
pids & fork bomb 填满进程表 & pids.current/max、fork 失败 \\
I/O & 大量写回拖慢共享磁盘 & io latency、队列深度、权重 \\
devices & 访问宿主敏感设备 & 设备拒绝日志和 errno \\
\bottomrule
\end{longtable}

资源限制要和观测一起配置。只设置 memory.max 但不记录 OOM 事件，故障时只能看到进程消失；只设置 CPU quota 但不看 throttled time，容易误判成应用锁或网络慢。

\subsection{seccomp：收缩系统调用面}

seccomp 允许进程限制自己可调用的系统调用。最小化系统调用面能减少内核攻击入口。网络服务正常处理请求可能只需要 \code{read/write/recv/send/epoll/futex/clock/mmap/brk/close} 等有限集合，不应该在运行阶段调用 \code{mount}、\code{ptrace}、\code{init\_module}。

\begin{lstlisting}
seccomp_policy:
  allow read, write, close
  allow epoll_wait, epoll_ctl
  allow recvfrom, sendto, accept4
  allow futex, clock_gettime
  allow mmap, munmap, brk
  deny mount, ptrace, bpf, init_module
\end{lstlisting}

allowlist 要从实际 trace 推导，而不是凭感觉写。测试时先在日志模式收集正常路径系统调用，再切到拒绝模式；同时要测试危险调用确实失败。seccomp 失败动作也要明确：返回 errno、kill 线程、kill 进程或记录审计。

\subsection{安全日志：可审计但不泄密}

安全事件日志需要足够结构化，便于按主体、对象、操作和原因查询。日志字段应包含 request id、pid、uid、容器/cgroup、操作、对象标识、策略名称和结果。敏感数据如 token、密码、密钥、完整请求体不应写入普通日志。

\begin{lstlisting}
security_event:
  request_id=...
  pid=...
  uid=...
  cgroup=...
  operation=open_write
  object=/data/report
  policy=readonly_mount
  result=denied
  errno=EROFS
\end{lstlisting}

日志也要有完整性和留存策略。攻击者若能删除或篡改日志，事后审计会失效；日志若无限增长，又会变成磁盘耗尽风险。高安全系统通常把安全事件发送到独立收集端，并给本地日志设置限额和轮转。

\subsection{性能观测：USE 与 RED 方法}

系统性能观测可以用两套简单框架。USE 关注资源：Utilization、Saturation、Errors，即资源使用率、排队饱和度和错误。RED 关注服务请求：Rate、Errors、Duration，即请求速率、错误率和耗时分布。

\begin{longtable}{p{0.18\linewidth}p{0.34\linewidth}p{0.36\linewidth}}
\toprule
框架 & 适用对象 & 指标例子 \\
\midrule
USE & CPU、内存、磁盘、网卡、锁、队列 & CPU busy、runqueue、EIO、drop、reclaim \\
RED & HTTP/RPC/数据库请求 & QPS、5xx、p50/p99、超时 \\
\bottomrule
\end{longtable}

这两套方法要结合。用户看到的是 RED：请求慢或失败；工程师定位要落到 USE：哪个资源饱和或错误。比如 p99 上升可能来自 CPU runnable wait、锁等待、page fault、dirty writeback、网络重传或 cgroup throttling。

\subsection{tracepoint、采样和事件日志}

观测手段有不同成本。计数器成本低，但细节少；事件 trace 细节多，但可能扰动系统；采样 profiler 能看到热点，但对短暂事件和 off-CPU 等待不敏感；分布式 trace 能串请求路径，但需要上下文传播。

\begin{longtable}{p{0.2\linewidth}p{0.34\linewidth}p{0.34\linewidth}}
\toprule
手段 & 适合回答 & 局限 \\
\midrule
计数器 & 是否发生、发生多少次 & 难以解释单个请求 \\
直方图 & 延迟分布和尾部 & 需要合理 bucket \\
事件 trace & 谁在何时做了什么 & 数据量大，需过滤 \\
CPU sampling & on-CPU 热点在哪里 & 看不到 sleeping 时间 \\
off-CPU trace & 线程等待在哪里 & 栈采集和符号成本 \\
\bottomrule
\end{longtable}

一个成熟诊断通常先看低成本指标定位方向，再打开更细粒度 trace 验证假设。直接全量抓所有事件，往往会制造过多噪声和额外开销。

\subsection{panic、oops 和可恢复错误}

内核错误并非都要 panic。用户传坏指针应返回 \code{EFAULT}；设备 I/O 失败应向上返回 \code{EIO}；某个驱动检测到可重置设备错误，可以隔离设备并恢复。panic 适用于内核不变量被破坏且继续运行可能造成更大损坏的情况，例如内核栈破坏、关键锁结构损坏、无法恢复的内存一致性错误。

\begin{longtable}{p{0.22\linewidth}p{0.36\linewidth}p{0.3\linewidth}}
\toprule
错误类型 & 处理 & 证据 \\
\midrule
用户错误 & 返回 errno 或发信号 & syscall、pid、地址、errno \\
设备错误 & 重试、reset、上报 EIO & 设备状态、请求 id、超时 \\
资源不足 & ENOMEM/OOM/限流 & cgroup、分配大小、victim \\
内核 bug & oops 或 panic & 寄存器、栈、锁、CPU、taint \\
\bottomrule
\end{longtable}

panic 路径要少依赖。此时文件系统、调度器或锁可能已经不可信，所以 crash dump 和 console 输出应尽量使用简单、预留或早期初始化的路径。复杂日志系统在 panic 中可能再次阻塞或破坏证据。

\subsection{故障注入：证明错误路径真的可用}

错误路径如果不测试，基本等于没实现。故障注入让内核或应用在可控点返回错误：分配失败、copy 用户失败、磁盘第 N 次写失败、网络短写、时钟跳变、锁超时、设备 completion 丢失。

\begin{lstlisting}
fault_injection_plan:
  fail allocation at callsite X once
  force copy_from_user to fail after 128 bytes
  return EIO on block write number N
  drop network response after commit
  delay wakeup by 500ms
\end{lstlisting}

每个注入点都要有预期：资源是否释放，引用计数是否回到原值，目标对象是否保持不变，错误是否传给调用者，日志是否足够诊断。故障注入把“理论上能处理错误”变成可验证事实。

\subsection{安全与观测的最终验收}

\begin{rubricbox}
一个边界合格的系统必须同时满足三件事：正常路径能完成；越权路径明确失败且没有副作用；失败路径留下足够证据，能说明主体、客体、操作、原因和影响范围。
\end{rubricbox}
